%% Question bank chapter file - generated by bank_tools.py build (do not edit while the bank is still being built:
%% it is regenerated from the per-test data in Question Bank\data). Every question is one bankq block:
%%   \begin{bankq}[id=..., source={...}, paper=.., qno=.., type=..., subtopic={...}, answer={...}, answer_src=key|solved,
%%                 working={...}, difficulty=1-10, diff_parts={K S I C F}, diff_note={...},
%%                 group=<passage id or empty>, also_in={duplicate copies in other tests}]
%% difficulty = min(10, K+S+I+C+F-1): K concepts 1-3, S steps 1-4, I insight 0-2, C computation 0-2, F format 0-1
%%   ...question LaTeX, options, figures (\bfig / \bfigi -> img/)...
%%   \end{bankq}
%% Passages shared by several questions are bankpassage blocks; their questions carry group=<passage id>.
\documentclass[10pt]{article}
\usepackage[a4paper,margin=1.8cm]{geometry}
\usepackage[T1]{fontenc}
\usepackage{amsmath,amssymb,mathtools}
\usepackage{graphicx}
\usepackage[export]{adjustbox}
\usepackage{xcolor,needspace,keyval}
\usepackage{enumitem}
\usepackage{booktabs,array,tabularx}
\usepackage[hidelinks]{hyperref}
\setlength{\parindent}{0pt}
\setlength{\parskip}{3pt}
\makeatletter
\define@key{bq}{id}{\def\bq@id{#1}}\define@key{bq}{source}{\def\bq@src{#1}}\define@key{bq}{paper}{\def\bq@paper{#1}}
\define@key{bq}{qno}{\def\bq@qno{#1}}\define@key{bq}{type}{\def\bq@type{#1}}\define@key{bq}{subtopic}{\def\bq@sub{#1}}
\define@key{bq}{answer}{\def\bq@ans{#1}}\define@key{bq}{answer_src}{\def\bq@asrc{#1}}\define@key{bq}{working}{\def\bq@work{#1}}
\define@key{bq}{group}{\def\bq@group{#1}}\define@key{bq}{also_in}{\def\bq@also{#1}}
\define@key{bq}{difficulty}{\def\bq@diff{#1}}\define@key{bq}{diff_parts}{\def\bq@dparts{#1}}\define@key{bq}{diff_note}{\def\bq@dnote{#1}}
\newenvironment{bankq}[1][]{\def\bq@work{}\def\bq@also{}\def\bq@group{}\def\bq@diff{?}\def\bq@dparts{}\def\bq@dnote{}\setkeys{bq}{#1}%
  \par\needspace{4\baselineskip}\medskip\noindent{\small\textbf{\bq@id}\hfill\color{gray}[\bq@src, Paper \bq@paper, Q\bq@qno]
  \bq@type\ $\cdot$ \bq@sub\ $\cdot$ \textbf{D\bq@diff}}\par\nopagebreak\relax}
  {\par{\small\color{teal}Answer: \textbf{\bq@ans}\ifx\bq@work\@empty\else\ (solved: \bq@work)\fi
  \ifx\bq@also\@empty\else\par\color{gray}Also in: \bq@also\fi}\par}
\define@key{bp}{id}{\def\bp@id{#1}}\define@key{bp}{source}{\def\bp@src{#1}}\define@key{bp}{questions}{\def\bp@qs{#1}}
\newenvironment{bankpassage}[1][]{\setkeys{bp}{#1}\par\needspace{4\baselineskip}\bigskip\noindent
  {\small\textbf{Passage \bp@id}\hfill\color{gray}[\bp@src; questions \bp@qs]}\par\nopagebreak\relax\itshape}{\par}
\makeatother
\newcommand{\bfig}[2]{\par\smallskip{\centering\includegraphics[width=#1pt,max width=0.62\linewidth,max height=0.32\textheight]{img/#2}\par}\smallskip}
\newcommand{\bfigi}[2]{\includegraphics[width=#1pt,max width=\linewidth,max height=0.32\textheight,valign=c]{img/#2}}
\newcommand{\bmissing}[1]{\textit{[figure missing in the source conversion: \detokenize{#1}]}}
\newenvironment{qoptions}{\begin{itemize}[leftmargin=1.8em,itemsep=1pt,topsep=2pt,label={}]}{\end{itemize}}
\newcommand{\opt}[2]{\item[(#1)] #2}
\providecommand{\fwstart}{}\providecommand{\fwend}{}\providecommand{\qfullsplit}{}
\newcolumntype{L}[1]{>{\raggedright\arraybackslash}p{#1}}
\newcolumntype{C}[1]{>{\centering\arraybackslash}p{#1}}
\newcolumntype{R}[1]{>{\raggedleft\arraybackslash}p{#1}}
\providecommand{\sectionhdr}[1]{}\providecommand{\subsectionhdr}[1]{}


\begin{document}
{\Large\bfseries SHM --- Question Bank}\par\smallskip
{\small Physics chapter P17. Questions from all converted test papers, grouped by source series. Answers: printed key where it matches the paper, otherwise solved (answer\_src=solved, with working). Exact repeats across papers appear once, with also\_in listing the other copies.}\par\medskip\hrule

\section*{Score-I papers}
\begin{bankq}[id=SI12-P1-PHY-S1-Q07, source={Score-I - Test 12 (09 Mar 2022)}, paper=1, qno=7, type=multi, subtopic={Spring-mass}, answer={A,C,D}, answer_src=key, difficulty=3, diff_parts={K1 S2 I0 C0 F1}, diff_note={Read amplitude, period and mean position off the x-t graph of SHM}]
A particle executes SHM about a point other than $x=0$ as shown in the graph. CHOOSE the correct option(s):
\bfig{137.3}{SI12_p1_phy_s1_q07_1.png}
\begin{qoptions}
\opt{A}{Amplitude is equal to 4 m}
\opt{B}{Angular frequency $=\dfrac\pi3$}
\opt{C}{Equilibrium position is at $x=2$ m}
\opt{D}{Angular frequency $=\dfrac{2\pi}3$}
\end{qoptions}
\end{bankq}

\section*{TOAS papers}
\begin{bankq}[id=TOAS01-P1-PHY-S1-Q01, source={TOAS - Test 01 (30 Jan 2022)}, paper=1, qno=1, type=single, subtopic={Spring-mass}, answer={A}, answer_src=key, difficulty=2, diff_parts={K1 S2 I0 C0 F0}, diff_note={F=m w^2 A and K=1/2 m w^2 A^2 solved for w}]
The maximum force on a body of mass $m$ executing SHM is $p$ and the maximum kinetic energy of the body is $q$. What is its time period?
\begin{qoptions}
\opt{A}{$\dfrac{2\pi}{p}\sqrt{2mq}$}
\opt{B}{$\dfrac{2\pi}{q}\sqrt{2mp}$}
\opt{C}{$\dfrac{\pi}{p}\sqrt{2mq}$}
\opt{D}{$\dfrac{\pi}{q}\sqrt{2mp}$}
\end{qoptions}
\end{bankq}
\begin{bankq}[id=TOAS01-P1-PHY-S1-Q02, source={TOAS - Test 01 (30 Jan 2022)}, paper=1, qno=2, type=single, subtopic={Energy}, answer={D}, answer_src=key, difficulty=1, diff_parts={K1 S1 I0 C0 F0}, diff_note={PE oscillates at twice the SHM frequency}]
Potential energy of a particle executing SHM varies sinusoidally. If the frequency of oscillation of the particle is $f$ then frequency of oscillation of potential energy is
\begin{qoptions}
\opt{A}{$f$}
\opt{B}{$\dfrac{f}{2}$}
\opt{C}{$4f$}
\opt{D}{$2f$}
\end{qoptions}
\end{bankq}
\begin{bankq}[id=TOAS01-P1-PHY-S2-Q05, source={TOAS - Test 01 (30 Jan 2022)}, paper=1, qno=5, type=multi, subtopic={Spring-mass}, answer={C,D}, answer_src=key, difficulty=3, diff_parts={K1 S2 I0 C0 F1}, diff_note={Four standard SHM statements; time T/12 vs T/6 is the only one needing a calculation}]
In a simple harmonic motion of amplitude $A$ (symbols have their usual meaning)
\begin{qoptions}
\opt{A}{Time taken by particle to move from mean position to $\dfrac{A}{2}$ is equal to time taken by particle to move from $\dfrac{A}{2}$ to $A$}
\opt{B}{Average velocity of particle in one time period $T$ is $\dfrac{4A}{T}$}
\opt{C}{Net work done on block in one time period $T$ is $0$}
\opt{D}{Particle has maximum kinetic energy at mean position}
\end{qoptions}
\end{bankq}
\begin{bankq}[id=TOAS01-P1-PHY-S2-Q06, source={TOAS - Test 01 (30 Jan 2022)}, paper=1, qno=6, type=multi, subtopic={Spring-mass}, answer={A}, answer_src=key, difficulty=5, diff_parts={K2 S3 I0 C0 F1}, diff_note={New equilibrium from k1 x1=k2 x2 plus k_eq=4000 gives T for each option}]
Two spring, each of outstretched length 20 cm but having different spring constants $k_1=1000$ N/m and $k_2=3000$ N/m are attached to two opposite faces of a block of mass $m=100$ g kept on a smooth horizontal surface as shown in figure.
\bfig{164.9}{TOAS01_p1_phy_s2_q06_1.png}
The outer ends of the two springs are now attached to two pins $P_1$ and $P_2$ whose locations are shown in the figure, as a result of this, the block acquired a new equilibrium position. The block has been displaced by small amount from its equilibrium position and released to perform simple harmonic motion, then.
\begin{qoptions}
\opt{A}{New equilibrium position is at 35 cm from $P_1$ and time period of simple harmonic motion is $\dfrac{\pi}{100}$ s.}
\opt{B}{New equilibrium position is at 20 cm from $P_1$ and time period of simple harmonic motion is $\dfrac{\pi}{100}$ s.}
\opt{C}{New equilibrium position is at 35 cm from $P_1$ and time period of simple harmonic motion is $\dfrac{\pi}{26}$ s.}
\opt{D}{New equilibrium position is at 30 cm from $P_1$ and time period of simple harmonic is $\dfrac{\pi}{26}$ s}
\end{qoptions}
\end{bankq}
\begin{bankq}[id=TOAS01-P1-PHY-S2-Q07, source={TOAS - Test 01 (30 Jan 2022)}, paper=1, qno=7, type=multi, subtopic={Energy}, answer={B,C,D}, answer_src=key, difficulty=5, diff_parts={K2 S2 I1 C0 F1}, diff_note={Differentiate v^2=a-bx^2 to get SHM, w^2=b and E at x=0}]
A particle of mass $m$ moves along X-axis. If $v$ is the velocity at a distance $x$ from a origin on the X-axis such that $v^2=a-bx^2$, where $a$ and $b$ are constant, then (Assume that potential energy of the particle is zero when it is at origin)
\begin{qoptions}
\opt{A}{the particle continues to move along the positive x-direction only}
\opt{B}{the motion is simple harmonic}
\opt{C}{the particle oscillates with a frequency equal to $\dfrac{\sqrt b}{2\pi}$}
\opt{D}{the total energy of the particle is $\dfrac{ma}{2}$}
\end{qoptions}
\end{bankq}
\begin{bankpassage}[id=TOAS01-P1-PHY-S3-G3-4, source={TOAS - Test 01 (30 Jan 2022)}, questions={3-4}]
A small object of mass 2 kg is placed on the x-axis in a potential field which potential energy is define as $U=16\,(x^2-2x)$ where position on x-axis, $x$ is in meter and potential energy $U$ is in joule. If the object is released from origin i.e. $x=0$. Then answer the following questions.
\end{bankpassage}
\begin{bankq}[id=TOAS01-P1-PHY-S3-Q03, source={TOAS - Test 01 (30 Jan 2022)}, paper=1, qno=3, type=numerical, subtopic={Energy}, answer={4}, answer_src=solved, working={$U^{\prime\prime}=32$ N/m, $m=2$ kg: $\omega=\sqrt{32/2}=4$ rad/s}, difficulty=2, diff_parts={K1 S1 I0 C0 F1}, diff_note={w=sqrt(U''/m) from the potential}, group=TOAS01-P1-PHY-S3-G3-4]
Angular frequency of oscillations ( in rad/second) is :
\end{bankq}
\begin{bankq}[id=TOAS01-P1-PHY-S3-Q04, source={TOAS - Test 01 (30 Jan 2022)}, paper=1, qno=4, type=numerical, subtopic={Energy}, answer={16}, answer_src=key, difficulty=3, diff_parts={K1 S2 I0 C0 F1}, diff_note={Equilibrium at x=1, released at x=0: KEmax = U(0)-U(1)}, group=TOAS01-P1-PHY-S3-G3-4]
Maximum kinetic energy (in joule) during oscillations is :
\end{bankq}
\begin{bankpassage}[id=TOAS01-P1-PHY-S3-G5-6, source={TOAS - Test 01 (30 Jan 2022)}, questions={5-6}]
A block $A$ of mass 1 kg connected with a spring is executing SHM. The displacement time equation of block is $x=8+4\sin(10t)$ (here $x$ is in m). An identical block $B$ moving towards negative x-axis with velocity 10 m/s collides with block $A$ and sticks to it at time $t=0$
\bfig{114.8}{TOAS01_p1_phy_s3_para5_6_1.png}
\end{bankpassage}
\begin{bankq}[id=TOAS01-P1-PHY-S3-Q05, source={TOAS - Test 01 (30 Jan 2022)}, paper=1, qno=5, type=numerical, subtopic={Spring-mass}, answer={15}, answer_src=solved, working={At $t=0$ block A is at its mean position: $v=4\times10=40$ m/s; momentum: $1\cdot40-1\cdot10=2v$, $v=15$ m/s}, difficulty=4, diff_parts={K2 S2 I0 C0 F1}, diff_note={A is at mean position at t=0 with v=40; perfectly inelastic collision with B}, group=TOAS01-P1-PHY-S3-G5-6]
Speed of block $A$ (in m/s) just after the collision is
\end{bankq}
\begin{bankq}[id=TOAS01-P1-PHY-S3-Q06, source={TOAS - Test 01 (30 Jan 2022)}, paper=1, qno=6, type=numerical, subtopic={Spring-mass}, answer={2.11 to 2.13}, answer_src=key, difficulty=5, diff_parts={K2 S3 I0 C0 F1}, diff_note={k=m w^2, new w=sqrt(k/2m), amplitude = v/w' at the mean position}, group=TOAS01-P1-PHY-S3-G5-6]
New amplitude (in m) of oscillation of block $A$ is
\end{bankq}
\begin{bankq}[id=TOAS01-P1-PHY-S4-Q07, source={TOAS - Test 01 (30 Jan 2022)}, paper=1, qno=7, type=numerical, subtopic={Spring-mass}, answer={25}, answer_src=key, difficulty=5, diff_parts={K2 S2 I1 C0 F1}, diff_note={Reduced mass of two blocks gives w; oscillations in 60 s}]
A spring of force constant $k=300$ N/m connects two blocks having masses 2 kg and 3 kg, lying on a smooth horizontal plane. If the spring block system is released from a stretched position, the number of complete oscillations in 1 minute is $6\times N$, find the value of $N$. Take $\pi=\sqrt{10}$.
\end{bankq}
\begin{bankq}[id=TOAS01-P1-PHY-S4-Q08, source={TOAS - Test 01 (30 Jan 2022)}, paper=1, qno=8, type=numerical, subtopic={Spring-mass}, answer={2}, answer_src=key, difficulty=3, diff_parts={K1 S2 I0 C0 F1}, diff_note={Amplitude 3mg/k about equilibrium: extension 4mg/k, compression 2mg/k}]
A block of mass $m$ hangs from a vertical spring of spring constant $k$. If the block is pulled down by a distance of $3\,mg/k$ from its equilibrium position and released, for the subsequent SHM of the block the ratio of maximum elongation in the spring to maximum compression in it is
\end{bankq}
\begin{bankq}[id=TOAS01-P1-PHY-S4-Q09, source={TOAS - Test 01 (30 Jan 2022)}, paper=1, qno=9, type=numerical, subtopic={Spring-mass}, answer={2}, answer_src=key, difficulty=6, diff_parts={K2 S3 I1 C0 F1}, diff_note={Moving pulley halves the speed; effective mass 5m/4 by energy, gravity cancels at equilibrium}]
If natural frequency $\omega$ of the system shown in figure is $x\sqrt{\dfrac{k}{5m}}$, find $x$.
\bfig{120.0}{TOAS01_p1_phy_s4_q09_1.png}
\end{bankq}
\begin{bankq}[id=TOAS01-P2-PHY-S1-Q02, source={TOAS - Test 01 (30 Jan 2022)}, paper=2, qno=2, type=single, subtopic={Spring-mass}, answer={A}, answer_src=key, difficulty=4, diff_parts={K2 S2 I1 C0 F0}, diff_note={Spring force component along the groove: k_eff=k cos^2(theta)}]
A block of mass $m$ is restricted to move in a smooth groove. It is connected to a spring of spring constant $k$. Neglecting the effect of gravity find time period of small oscillation :-
\bfig{103.6}{TOAS01_p2_phy_s1_q02_1.png}
\begin{qoptions}
\opt{A}{$\dfrac{2\pi}{\cos\theta}\sqrt{\dfrac mk}$}
\opt{B}{$2\pi\sqrt{\dfrac mk}$}
\opt{C}{$\dfrac{2\pi}{\sin\theta}\sqrt{\dfrac mk}$}
\opt{D}{$2\pi\sqrt{\dfrac{m\cos^2\theta}{k}}$}
\end{qoptions}
\end{bankq}
\begin{bankq}[id=TOAS01-P2-PHY-S1-Q03, source={TOAS - Test 01 (30 Jan 2022)}, paper=2, qno=3, type=single, subtopic={Spring-mass}, answer={C}, answer_src=key, difficulty=2, diff_parts={K1 S1 I1 C0 F0}, diff_note={Only the extreme position is visited once per period}]
The time period of a particle in simple harmonic motion is equal to the time between consecutive appearances of the particle at a particular point in its motion. This point is
\begin{qoptions}
\opt{A}{The mean position}
\opt{B}{Between the mean position and the positive extreme}
\opt{C}{An extreme position}
\opt{D}{Between the mean position and the negative extreme.}
\end{qoptions}
\end{bankq}
\begin{bankq}[id=TOAS01-P2-PHY-S1-Q05, source={TOAS - Test 01 (30 Jan 2022)}, paper=2, qno=5, type=single, subtopic={Energy}, answer={D}, answer_src=key, difficulty=2, diff_parts={K1 S2 I0 C0 F0}, diff_note={U''=U0 b^2 gives T}]
Potential energy associated with a particle of mass $m$ moving along $x$ axis is given by $U=U_0[1-\cos(bx)]$ where $U_0$ \& $b$ are positive constant. Find the time period of small oscillation
\begin{qoptions}
\opt{A}{$2\pi\sqrt{\dfrac{m}{U_0b}}$}
\opt{B}{$2\pi\sqrt{\dfrac{mb}{U_0}}$}
\opt{C}{$2\pi\sqrt{\dfrac{m}{U_0}}$}
\opt{D}{$\dfrac{2\pi}{b}\sqrt{\dfrac{m}{U_0}}$}
\end{qoptions}
\end{bankq}
\begin{bankq}[id=TOAS01-P2-PHY-S1-Q06, source={TOAS - Test 01 (30 Jan 2022)}, paper=2, qno=6, type=single, subtopic={Energy}, answer={D}, answer_src=key, difficulty=2, diff_parts={K1 S1 I1 C0 F0}, diff_note={E=kA^2/2 independent of the mass}]
The total mechanical energy of a spring-mass system in simple harmonic motion is $E=\dfrac12m\omega^2A^2$. Suppose the oscillating particle is replaced by another particle of double the mass while the amplitude $A$ remains the same. The new mechanical energy will
\begin{qoptions}
\opt{A}{become $2E$}
\opt{B}{become $E/2$}
\opt{C}{become $\sqrt2E$}
\opt{D}{remain $E$}
\end{qoptions}
\end{bankq}
\begin{bankq}[id=TOAS01-P2-PHY-S1-Q08, source={TOAS - Test 01 (30 Jan 2022)}, paper=2, qno=8, type=single, subtopic={Spring-mass}, answer={D}, answer_src=key, difficulty=3, diff_parts={K1 S2 I1 C0 F0}, diff_note={Half period with k plus half period with 2k}]
In the figure block lies on smooth horizontal surface and both the springs have same spring constant $k$. The spring $B$ is attached with block while spring $A$ just touches the block as the position shown and both the springs are in relaxed state. The time period of oscillation is :-
\bfig{94.2}{TOAS01_p2_phy_s1_q08_A.png}
\begin{qoptions}
\opt{A}{$2\pi\sqrt{\dfrac{m}{2k}}$}
\opt{B}{$2\pi\sqrt{\dfrac{3m}{2k}}$}
\opt{C}{$2\pi\sqrt{\dfrac mk}$}
\opt{D}{$\dfrac{(\sqrt2+1)}{\sqrt2}\pi\sqrt{\dfrac mk}$}
\end{qoptions}
\end{bankq}
\begin{bankq}[id=TOAS01-P2-PHY-S1-Q09, source={TOAS - Test 01 (30 Jan 2022)}, paper=2, qno=9, type=single, subtopic={Spring-mass}, answer={B}, answer_src=key, difficulty=1, diff_parts={K1 S1 I0 C0 F0}, diff_note={a=-w^2 A cos(wt) at wt=pi/4}]
A particle executes SHM about the point $x=0$. At time $t=0$ it has position $x=A$ and zero velocity. The angular frequency of oscillation is $\omega$. The acceleration of the particle at $t=\dfrac{\pi}{4\omega}$ is :-
\begin{qoptions}
\opt{A}{$-A\omega^2$}
\opt{B}{$-\dfrac{A\omega^2}{\sqrt2}$}
\opt{C}{Zero}
\opt{D}{$-A\omega^2\sqrt2$}
\end{qoptions}
\end{bankq}
\begin{bankq}[id=TOAS01-P2-PHY-S1-Q10, source={TOAS - Test 01 (30 Jan 2022)}, paper=2, qno=10, type=single, subtopic={Spring-mass}, answer={A}, answer_src=key, difficulty=4, diff_parts={K2 S2 I1 C0 F0}, diff_note={Rolling sphere effective mass 7m/5 with springs 2K, then x=x0 cos wt}]
Two springs each of spring constant $K$ are connected with the centre of a sphere of mass $m$ and radius $R$ as shown in figure. Given system lies on horizontal surface and both the springs are horizontal. Friction is sufficient to prevent relative slipping between horizontal surface and sphere. From the position where both springs are at natural length, sphere is slightly displaced by a distance `$x_0$' and released from rest. After being released from rest find minimum time taken by sphere to move by a distance $\dfrac{x_0}{2}$.
\bfig{105.1}{TOAS01_p2_phy_s1_q10_1.png}
\begin{qoptions}
\opt{A}{$\dfrac\pi3\sqrt{\dfrac{7m}{10k}}$}
\opt{B}{$\dfrac\pi3\sqrt{\dfrac{5m}{7k}}$}
\opt{C}{$\dfrac\pi4\sqrt{\dfrac{7m}{5k}}$}
\opt{D}{$\dfrac\pi6\sqrt{\dfrac{5m}{7k}}$}
\end{qoptions}
\end{bankq}
\begin{bankq}[id=TOAS01-P2-PHY-S1-Q11, source={TOAS - Test 01 (30 Jan 2022)}, paper=2, qno=11, type=single, subtopic={Pendulums}, answer={D}, answer_src=key, difficulty=4, diff_parts={K2 S2 I1 C0 F0}, diff_note={Pseudo-force gives g_eff=sqrt2 g for the pendulum}]
A simple pendulum is suspended from the ceiling of a car moving with constant speed on a horizontal circular track of radius $R$. Speed of the car is given by $v=\sqrt{gR}$ where $g$ is acceleration due to gravity and $R$ is radius of circular turn. Find the time period of a simple pendulum of length $\ell$ suspended from the ceiling of the car.
\begin{qoptions}
\opt{A}{$2\pi\sqrt{\dfrac{\ell}{2g}}$}
\opt{B}{$2\pi\sqrt{\dfrac{2\ell}{g}}$}
\opt{C}{$2\pi\sqrt{\dfrac{\ell}{g}}$}
\opt{D}{$2\pi\sqrt{\dfrac{\ell}{\sqrt2g}}$}
\end{qoptions}
\end{bankq}
\begin{bankq}[id=TOAS01-P2-PHY-S1-Q12, source={TOAS - Test 01 (30 Jan 2022)}, paper=2, qno=12, type=single, subtopic={Spring-mass}, answer={B}, answer_src=key, difficulty=2, diff_parts={K1 S2 I0 C0 F0}, diff_note={v=A w sin(wt)=A w/sqrt2 gives wt=pi/4}]
In a horizontal spring-mass system, at $t=0$, mass `$m$' is released from rest after being displaced (from NLP of spring) towards right by some distance on a frictionless surface. Find minimum time after which speed of block becomes $\dfrac{1}{\sqrt2}$ times of maximum possible speed during the motion.
\bfig{122.4}{TOAS01_p2_phy_s1_q12_A.png}
\begin{qoptions}
\opt{A}{$\dfrac\pi2\sqrt{\dfrac mk}$}
\opt{B}{$\dfrac\pi4\sqrt{\dfrac mk}$}
\opt{C}{$\pi\sqrt{\dfrac mk}$}
\opt{D}{$\dfrac\pi8\sqrt{\dfrac mk}$}
\end{qoptions}
\end{bankq}
\begin{bankq}[id=TOAS01-P2-PHY-S1-Q13, source={TOAS - Test 01 (30 Jan 2022)}, paper=2, qno=13, type=single, subtopic={Spring-mass}, answer={D}, answer_src=key, difficulty=4, diff_parts={K1 S3 I1 C0 F0}, diff_note={Phase of each particle from initial conditions, solve for crossing time}]
Two particles $A$ and $B$ are performing SHM along same straight line with same amplitude $a$ and same angular frequency $\omega$. At $t=0$, $A$ is at the left extreme and $B$ is at a distance $a/2$ from mean position to the right of mean position and moving toward mean position. Then:
\begin{qoptions}
\opt{A}{their phase difference is $\dfrac{5\pi}{3}$}
\opt{B}{their phase difference is $\dfrac{\pi}{3}$}
\opt{C}{time after which they cross each other is $\dfrac{2\pi}{3\omega}$}
\opt{D}{time after which they cross each other is $\dfrac{\pi}{3\omega}$}
\end{qoptions}
\end{bankq}
\begin{bankq}[id=TOAS01-P2-PHY-S1-Q14, source={TOAS - Test 01 (30 Jan 2022)}, paper=2, qno=14, type=single, subtopic={Energy}, answer={C}, answer_src=key, difficulty=4, diff_parts={K2 S3 I0 C0 F0}, diff_note={Read amplitude and k from U=kx^2/2 on the graph, KE max at equilibrium}]
Figure shows the potential-energy diagram and the total energy line of a particle oscillating on a spring. Which of the following option is incorrect.
\bfig{146.3}{TOAS01_p2_phy_s1_q14_1.png}
\begin{qoptions}
\opt{A}{the amplitude of oscillation of particle is 6 cm.}
\opt{B}{the spring constant is 4.44 kN/m.}
\opt{C}{the mass of particle is 1kg.}
\opt{D}{the maximum kinetic energy of particle is at $x=20$ cm.}
\end{qoptions}
\end{bankq}
\begin{bankq}[id=TOAS01-P2-PHY-S1-Q15, source={TOAS - Test 01 (30 Jan 2022)}, paper=2, qno=15, type=single, subtopic={Spring-mass}, answer={C}, answer_src=key, difficulty=6, diff_parts={K3 S3 I1 C0 F0}, diff_note={SHM max acceleration, incline component and friction limit on the upper block}]
The coefficient of friction between block of mass $m$ and $2m$ is $\mu=2\tan\theta$. There is no friction between block of mass $2m$ and inclined plane. The maximum amplitude of two block system for which there is no relative motion between both the blocks is:
\bfig{97.3}{TOAS01_p2_phy_s1_q15_1.png}
\begin{qoptions}
\opt{A}{$g\sin\theta\sqrt{\dfrac km}$}
\opt{B}{$\dfrac{mg\sin\theta}{k}$}
\opt{C}{$\dfrac{3mg\sin\theta}{k}$}
\opt{D}{None of these}
\end{qoptions}
\end{bankq}
\begin{bankq}[id=TOAS01-P2-PHY-S1-Q17, source={TOAS - Test 01 (30 Jan 2022)}, paper=2, qno=17, type=single, subtopic={Pendulums}, answer={C}, answer_src=key, difficulty=6, diff_parts={K2 S3 I1 C1 F0}, diff_note={Physical pendulum T^2 ~ (L^2/12+d^2)/d minimised at d=L/sqrt12}]
A rod of mass `$m$' and length `$L$' is suspended in vertical plane so as to oscillate freely about an axis normal to the plane of paper passing through `$S$' (point of suspension). Find the distance of `$S$' from end $A$ of rod such that time period of small oscillation is minimum.
\bfig{22.0}{TOAS01_p2_phy_s1_q17_1.png}
\begin{qoptions}
\opt{A}{$\dfrac L2$}
\opt{B}{$\dfrac{L}{\sqrt{12}}$}
\opt{C}{$\dfrac L2+\dfrac{L}{\sqrt{12}}$}
\opt{D}{$\dfrac L2-\dfrac{L}{\sqrt{12}}$}
\end{qoptions}
\end{bankq}
\begin{bankq}[id=TOAS01-P2-PHY-S1-Q19, source={TOAS - Test 01 (30 Jan 2022)}, paper=2, qno=19, type=single, subtopic={Spring-mass}, answer={D}, answer_src=key, difficulty=4, diff_parts={K2 S2 I1 C0 F0}, diff_note={No separation needs maximum acceleration g, w^2=k/(M+m)}]
A small block of mass $m$ is kept on a bigger block of mass $M$ which is attached to a vertical spring of spring constant $k$ as shown in the figure. The system oscillates vertically. What can be the maximum amplitude with which the two blocks may oscillate together ?
\bfig{52.7}{TOAS01_p2_phy_s1_q19_1.png}
\begin{qoptions}
\opt{A}{$\dfrac{2(M+m)g}{k}$}
\opt{B}{$\dfrac{(m)g}{k}$}
\opt{C}{$\dfrac{(M)g}{k}$}
\opt{D}{$\dfrac{(M+m)g}{k}$}
\end{qoptions}
\end{bankq}
\begin{bankq}[id=TOAS01-P2-PHY-S2-Q03, source={TOAS - Test 01 (30 Jan 2022)}, paper=2, qno=3, type=numerical, subtopic={Spring-mass}, answer={0.20}, answer_src=key, difficulty=3, diff_parts={K1 S2 I0 C0 F1}, diff_note={T=1/1000 s, v_max=A w=200 m/s}]
Equation of motion of a particle performing SHM is given by\\
$y=\dfrac{10}{\pi}\sin\left[2000\pi t-\dfrac\pi6\right]$, where $y$ is in cm and $t$ is in sec.\\
If time period of SHM is $t$ and maximum possible velocity of the particle during SHM is $v_{\max}$ then find the value of $t\times v_{\max}$ (in SI unit).
\end{bankq}
\begin{bankq}[id=TOAS01-P2-PHY-S2-Q04, source={TOAS - Test 01 (30 Jan 2022)}, paper=2, qno=4, type=numerical, subtopic={Pendulums}, answer={0.25}, answer_src=key, difficulty=4, diff_parts={K2 S2 I0 C0 F1}, diff_note={Torsional pendulum T=2pi sqrt(I/kappa), I=mR^2/2}]
A uniform disc of radius 5.0 cm and mass 200g is fixed at its centre to a metal wire, the other end of which is fixed with a clamp. The hanging disc is rotated about the wire through an angle and a released. If the disc makes torsional oscillations with time period 0.20 s, find the torsional constant (in S.I. units) of the wire. (Take : $\pi^2=10$)
\bfig{34.6}{TOAS01_p2_phy_s2_q04_1.png}
\end{bankq}
\begin{bankq}[id=TOAS01-P2-PHY-S2-Q08, source={TOAS - Test 01 (30 Jan 2022)}, paper=2, qno=8, type=numerical, subtopic={Pendulums}, answer={15}, answer_src=key, difficulty=4, diff_parts={K2 S2 I0 C0 F1}, diff_note={Buoyancy gives g_eff=15g/16 for the pendulum}]
A simple pendulum oscillating in air has period $T$. The bob of the pendulum is completely immersed in a non-viscous liquid. The density of the liquid is $\dfrac1{16}$ th of the density of material of the bob. If bob is inside liquid all the time and its period of oscillation in this liquid is given by $4T\sqrt{\dfrac1N}$ then find the value of $N$.
\end{bankq}
\begin{bankq}[id=TOAS01-P2-PHY-S2-Q10, source={TOAS - Test 01 (30 Jan 2022)}, paper=2, qno=10, type=numerical, subtopic={Pendulums}, answer={3}, answer_src=key, difficulty=5, diff_parts={K2 S3 I0 C0 F1}, diff_note={I=7mL^2/12 about the wire, energy conservation gives v=theta0 sqrt(3k/7m)}]
Two small balls each of mass $m$ are connected by a uniform rod of mass $m$ and length $L$. The given system is suspended from its centre by a thin wire of torsional constant $k$. The rod is rotated about the wire through an angle $\theta_0$ and released. Find the velocity of each ball as the system passes through the mean position. If this velocity is given by $v=\theta_0\sqrt{\dfrac{\alpha k}{7m}}$ then fill the value of $\alpha$.
\bfig{108.4}{TOAS01_p2_phy_s2_q10_1.png}
\end{bankq}
\begin{bankq}[id=TOAS02-P1-PHY-S1-Q02, source={TOAS - Test 02 (13 Feb 2022)}, paper=1, qno=2, type=single, subtopic={Spring-mass}, answer={C}, answer_src=key, difficulty=4, diff_parts={K1 S3 I1 C0 F0}, diff_note={Displacement over a quarter period from x=A sin(wt+theta), divided by T/4}]
The displacement of particle executing SHM is given by $x=A\sin(\omega t+\theta)$, where $\theta$-initial phase of motion, $A$-Amplitude of motion and $T$-Time period. The magnitude of average velocity during motion of particle between $t=0$ to $t=T/4$ is :
\begin{qoptions}
\opt{A}{$2\dfrac{A\omega}{\pi}$}
\opt{B}{$2\dfrac{A\omega}{\pi}|\cos\theta|$}
\opt{C}{$2\dfrac{A\omega}{\pi}|\cos\theta-\sin\theta|$}
\opt{D}{$2\dfrac{A\omega}{\pi}|\cos\theta+\sin\theta|$}
\end{qoptions}
\end{bankq}
\begin{bankq}[id=TOAS02-P1-PHY-S1-Q03, source={TOAS - Test 02 (13 Feb 2022)}, paper=1, qno=3, type=single, subtopic={Spring-mass}, answer={A}, answer_src=key, difficulty=5, diff_parts={K2 S2 I1 C1 F0}, diff_note={Two perpendicular SHMs: d^2=29-20cos(wt-phi) via R-cosine form, d from 3 to 7}]
Two particles $A$ and $B$ are performing SHM along $x$ and $y$-axis respectively with equal amplitude and frequency of 2 cm and 1 Hz respectively. Equilibrium positions of the particles $A$ and $B$ are at the co-ordinates (3, 0) cm and (0, 4) cm respectively. At $t=0$, $B$ is at its equilibrium position and moving towards the origin, while $A$ is nearest to the origin and moving away from the origin. If smallest separation between $A$ and $B$ is $d_{\min}$ and maximum separation is $d_{\max}$ then $d_{\max}-d_{\min}=$ \underline{\hspace{2.5em}} ?
\begin{qoptions}
\opt{A}{4 cm}
\opt{B}{5 cm}
\opt{C}{2 cm}
\opt{D}{1.5 cm}
\end{qoptions}
\end{bankq}
\begin{bankq}[id=TOAS02-P1-PHY-S2-Q12, source={TOAS - Test 02 (13 Feb 2022)}, paper=1, qno=12, type=multi, subtopic={Spring-mass}, answer={B,C}, answer_src=key, difficulty=6, diff_parts={K1 S3 I1 C1 F1}, diff_note={x=a(1-cos(2wt-pi/2))/2 gives amplitude a/2... then v^2 relation from energy form}]
A particle moves along the x-axis according to the law $x=a\sin^2\left(\omega t-\dfrac\pi4\right)$
\begin{qoptions}
\opt{A}{The amplitude of the oscillation is $a$}
\opt{B}{The time period of oscillations is $\dfrac\pi\omega$}
\opt{C}{The velocity ($v$) of the particle is given by $v^2=4\omega^2x(a-x)$}
\opt{D}{The velocity ($v$) of the particle is given by $v^2=2\omega^2x(a-x)$}
\end{qoptions}
\end{bankq}
\begin{bankq}[id=TOAS02-P1-PHY-S3-Q05, source={TOAS - Test 02 (13 Feb 2022)}, paper=1, qno=5, type=numerical, subtopic={Spring-mass}, answer={500}, answer_src=key, difficulty=6, diff_parts={K2 S3 I1 C0 F1}, diff_note={Rotational SHM with torsional constant k r^2 from the spring, period read from the graph}]
A metal plate is mounted on an axle through its centre of mass. A spring with spring constant $k=2000$ N/m connects a wall with a point on the rim a distance $r=2.5$ cm from the centre of mass. Initially the spring is at its rest length. If the plate is rotated by $7^\circ$ and released, it rotates about the axle in SHM with its angular position given by the graph shown. The horizontal axis scale is set by $t_s=31.4$ ms. If the moment of inertia of the plate about the centre of mass is $x\times10^{-6}\ \mathrm{kg\,m^2}$. Then the value of $x$ is:
\bfig{245.2}{TOAS02_p1_phy_s3_q05_1.png}
\end{bankq}
\begin{bankq}[id=TOAS03-P1-PHY-S1-Q02, source={TOAS - Test 03 (27 Feb 2022)}, paper=1, qno=2, type=single, subtopic={Energy}, answer={C}, answer_src=key, difficulty=3, diff_parts={K1 S2 I1 C0 F0}, diff_note={PE is proportional to displacement squared so sqrt(E)=sqrt(E1)+sqrt(E2)}]
A body is executing simple harmonic motion. At a displacement $x$, its potential energy is $E_1$ and at a displacement of $y$, its potential energy is $E_2$. The potential energy $E$ at a displacement $(x+y)$ is :-
\begin{qoptions}
\opt{A}{$E_1+E_2$}
\opt{B}{$\sqrt{E_1^2+E_2^2}$}
\opt{C}{$E_1+E_2+2\sqrt{E_1E_2}$}
\opt{D}{$\sqrt{E_1E_2}$}
\end{qoptions}
\end{bankq}
\begin{bankq}[id=TOAS03-P1-PHY-S1-Q04, source={TOAS - Test 03 (27 Feb 2022)}, paper=1, qno=4, type=single, subtopic={Spring-mass}, answer={A}, answer_src=key, difficulty=5, diff_parts={K2 S3 I1 C0 F0}, diff_note={Spring-coupled blocks that separate: motion until Q reaches its lowest point and the gap}]
In the figure shown, blocks $P$ and $Q$ are in contact but do not stick to each other. The springs are in their natural lengths. The system is released from rest. Then the distance between $P$ and $Q$ when $Q$ is at the lowest point for first time will be :-
\bfig{45.6}{TOAS03_p1_phy_s1_q04_1.png}
\begin{qoptions}
\opt{A}{$\dfrac{2mg}{K}$}
\opt{B}{$\dfrac{4mg}{K}$}
\opt{C}{$\dfrac{3mg}{K}$}
\opt{D}{0}
\end{qoptions}
\end{bankq}
\begin{bankq}[id=TOAS03-P1-PHY-S2-Q08, source={TOAS - Test 03 (27 Feb 2022)}, paper=1, qno=8, type=multi, subtopic={Energy}, answer={A,D}, answer_src=key, difficulty=4, diff_parts={K1 S3 I0 C0 F1}, diff_note={KE-time graph: KE max and period read, amplitude from KE_max=1/2 m w^2 A^2}]
Kinetic energy of a particle executing SHM changes with time according to following graph. Mass of the particle is 2 kg. Then :-
\bfig{210.6}{TOAS03_p1_phy_s2_q08_1.png}
\begin{qoptions}
\opt{A}{amplitude of oscillation is 0.4 m}
\opt{B}{acceleration of the particle is zero at $t=0.1$ s}
\opt{C}{time period of oscillation is 0.1 s}
\opt{D}{frequency of oscillation is 5 Hz.}
\end{qoptions}
\end{bankq}
\begin{bankq}[id=TOAS03-P1-PHY-S2-Q10, source={TOAS - Test 03 (27 Feb 2022)}, paper=1, qno=10, type=multi, subtopic={Energy}, answer={A,B,C,D}, answer_src=key, difficulty=5, diff_parts={K2 S3 I0 C0 F1}, diff_note={U=2y^2-30y+450: equilibrium 7.5, w=sqrt(k/m), amplitude from energy, four options}]
A particle of mass $m=3$ kg is moving on a straight line in a conservative force field in which the potential energy associated with the particle is given by $U(y)=2y^2-30y+450$. At $t=0$, particle is at $y=8$ m and moving towards positive $y$-axis with 1 m/s.
\begin{qoptions}
\opt{A}{During the course of its motion the maximum speed of the particle will be $\sqrt{\dfrac43}$ m/s}
\opt{B}{Particle will never cross the origin.}
\opt{C}{Particle will come to its equilibrium position for the first time at $\dfrac{5\pi}{4\sqrt3}$ second.}
\opt{D}{The amplitude of oscillation for the particle is 1 m.}
\end{qoptions}
\end{bankq}
\begin{bankq}[id=TOAS03-P1-PHY-S3-Q02, source={TOAS - Test 03 (27 Feb 2022)}, paper=1, qno=2, type=numerical, subtopic={Pendulums}, answer={4.00}, answer_src=key, difficulty=4, diff_parts={K2 S2 I0 C0 F1}, diff_note={I of a hemisphere about its axis (2/5 m R^2) and T=2 pi sqrt(I/C)}]
A uniform solid hemisphere of mass $m$ and radius $R$ is attached to the roof with a chord of torsional constant $C$ and performing torsional SHM. Then the time period (in seconds) of SHM is (Take $m=15$ kg, $R=\dfrac2\pi$ m, $C=6$ Nm/rad.)
\bfig{66.4}{TOAS03_p1_phy_s3_q02_1.png}
\end{bankq}
\begin{bankq}[id=TOAS03-P1-PHY-S3-Q04, source={TOAS - Test 03 (27 Feb 2022)}, paper=1, qno=4, type=numerical, subtopic={Spring-mass}, answer={2}, answer_src=solved, working={Both particles have the same velocity 1 m/s at positions $\pm x_0$, so phases $\pm\theta$: $A\omega\cos\theta=1$; they coincide when $\omega t=\pi/2$ where the speed is $A\omega\sin\theta=\sqrt3$; so $(A\omega)^2=1+3$ and $v_{max}=A\omega=2$ m/s}, difficulty=5, diff_parts={K1 S3 I1 C0 F1}, diff_note={Same amplitude and equal velocities at +-x0: phases +-theta; speed sqrt3 where they meet gives A w}]
Two particles are performing SHM with same amplitude and time period. At an instant two particles are having velocity 1 m/s but one is on the right and the other is on the left of their mean position. When the particles have same position, their speed is $\sqrt3$ m/s. Find the maximum speed (in m/s) of particles during SHM.
\end{bankq}
\begin{bankq}[id=TOAS03-P2-PHY-S1-Q01, source={TOAS - Test 03 (27 Feb 2022)}, paper=2, qno=1, type=multi, subtopic={Spring-mass}, answer={A,B,D}, answer_src=key, difficulty=7, diff_parts={K2 S4 I0 C1 F1}, diff_note={Energy conservation gives max speed and maximum compression; time from phase for four options}]
A small block of mass $m$ is fixed at upper end of a massless vertical spring of spring constant $K=\dfrac{4mg}{L}$ and natural length 10 m. The lower end of spring is free and is at a height $L$ meter from fixed horizontal floor as shown. The spring is initially unstretched and the spring block system is released from rest in the shown position.
\bfig{92.5}{TOAS03_p2_phy_s1_q01_1.png}
\begin{qoptions}
\opt{A}{When the speed of block is maximum the magnitude of force exerted by spring on the block is $mg$.}
\opt{B}{The maximum speed attained by the block is $\dfrac32\sqrt{gL}$.}
\opt{C}{The maximum compression of the spring is $\dfrac L8$}
\opt{D}{The time taken by the mass to move from the equilibrium position to the maximum compressed position of the spring is $\dfrac\pi4\sqrt{\dfrac Lg}$}
\end{qoptions}
\end{bankq}
\begin{bankq}[id=TOAS03-P2-PHY-S2-Q05, source={TOAS - Test 03 (27 Feb 2022)}, paper=2, qno=5, type=numerical, subtopic={Spring-mass}, answer={2.50}, answer_src=key, difficulty=6, diff_parts={K2 S3 I1 C0 F1}, diff_note={Combine series and parallel springs and effective mass for the oscillation period}]
In the figure shown below all the springs are ideal and time period of small oscillations about equilibrium position is $\pi\sqrt{\dfrac{xm}{k}}$. Find the value of $x$.
\bfig{174.2}{TOAS03_p2_phy_s2_q05_1.png}
\end{bankq}
\begin{bankq}[id=TOAS03-P2-PHY-S3-Q05, source={TOAS - Test 03 (27 Feb 2022)}, paper=2, qno=5, type=numerical, subtopic={Energy}, answer={3}, answer_src=key, difficulty=4, diff_parts={K1 S3 I0 C0 F1}, diff_note={U''=360 gives w=6 sqrt10 and T=1/n with pi^2=10}]
The potential energy of a particle of mass 1 kg in motion along the $x$-axis is given by $U=(10-10\cos6x)$ J. If the period of small oscillations is $\dfrac1n$ second, find the value of $n$. $\left(\pi^2\approx10\right)$.
\end{bankq}
\begin{bankq}[id=TOAS04-P1-PHY-S2-Q07, source={TOAS - Test 04 (15 Mar 2022)}, paper=1, qno=7, type=multi, subtopic={Spring-mass}, answer={A,B,C}, answer_src=key, difficulty=8, diff_parts={K2 S3 I2 C1 F1}, diff_note={Pulse impulse and w=pi: pulses at quarter periods, amplitudes combine as sqrt(0.1^2+0.1^2)}]
A block of mass 1 gm carrying a charge $1\,\mu$C is connected to a spring of stiffness $(\pi^2\times10^{-3})$ N/m whose other end is fixed to a wall. The block rests on a smooth ground as shown. By making a suitable arrangement, an electric field lying in the plane of the figure, directed towards right with strength $\left(\dfrac\pi{10}\right)$ MV/m is switched on exactly for an interval of 1 ms after every half second. Take the moment when the field is switched on for first time as $t=0$. Now, select the correct options.
\bfig{132.7}{TOAS04_p1_phy_s2_q07_1.png}
\begin{qoptions}
\opt{A}{The maximum extension in the spring is 0.1414 m (nearly)}
\opt{B}{The spring never gets compressed.}
\opt{C}{The time taken by the block to start from mean position and reach the position of maximum elongation of spring is 0.75s.}
\opt{D}{The time taken by the block to start from mean position to the position of maximum elongation of spring is 0.5 s.}
\end{qoptions}
\end{bankq}
\begin{bankpassage}[id=TOAS04-P1-PHY-S3-G3-4, source={TOAS - Test 04 (15 Mar 2022)}, questions={3-4}]
In the given figure, the block is attached with a system of three ideal springs $A$, $B$, $C$. The block is displaced slightly by a distance `$x$' from its equilibrium position vertically downwards and released.
\bfig{40.1}{TOAS04_p1_phy_s3_para3_4_1.png}
\end{bankpassage}
\begin{bankq}[id=TOAS04-P1-PHY-S3-Q03, source={TOAS - Test 04 (15 Mar 2022)}, paper=1, qno=3, type=numerical, subtopic={Spring-mass}, answer={4.00}, answer_src=key, difficulty=4, diff_parts={K1 S2 I1 C0 F1}, diff_note={Series-parallel springs: ratio of deformations from the force split}, group=TOAS04-P1-PHY-S3-G3-4]
The ratio of deformation of spring $B$ to spring $C$ is
\end{bankq}
\begin{bankq}[id=TOAS04-P1-PHY-S3-Q04, source={TOAS - Test 04 (15 Mar 2022)}, paper=1, qno=4, type=numerical, subtopic={Spring-mass}, answer={5.50}, answer_src=key, difficulty=6, diff_parts={K2 S3 I1 C0 F1}, diff_note={Effective spring constant of the three-spring network and the period}, group=TOAS04-P1-PHY-S3-G3-4]
The time period of small vertical oscillations of block is $2\pi\sqrt{a\dfrac mk}$. Find $a$
\end{bankq}
\begin{bankq}[id=TOAS05-P1-PHY-S2-Q09, source={TOAS - Test 05 (03 Apr 2022)}, paper=1, qno=9, type=single, subtopic={Spring-mass}, answer={A}, answer_src=key, difficulty=4, diff_parts={K1 S3 I1 C0 F0}, diff_note={SHM with phi=0: times for rest, max velocity, KE=PE and half max velocity}]
A particle is performing SHM $x=A\sin(\omega t+\phi)$ whose time period is $T$. List-I contains events related to the motion of particle and List-II contains the minimum time after which that event happen. Choose the correct match from option if $\phi=0^\circ$
{\small\setlength{\tabcolsep}{3pt}\renewcommand{\arraystretch}{1.25}\begin{center}\begin{tabular}{|L{0.75cm}|L{3.6cm}|L{0.7cm}|L{1.0cm}|}
\hline
\multicolumn{2}{|c|}{List--I} & \multicolumn{2}{c|}{List--II}\\
\hline
(I) & Particle comes to rest & (P) & $T/4$\\
\hline
(II) & Velocity is maximum & (Q) & $T/6$\\
\hline
(III) & Kinetic energy = Potential energy & (R) & $3T/8$\\
\hline
(IV) & Velocity is half of maximum velocity & (S) & $T/24$\\
\hline
\multicolumn{2}{|c|}{} & (T) & $T/2$\\
\hline
\multicolumn{2}{|c|}{} & (U) & $T/8$\\
\hline
\end{tabular}\end{center}}
\begin{qoptions}
\opt{A}{I $\to$ P;II $\to$ T;III $\to$ U;IV $\to$ Q}
\opt{B}{I $\to$ Q;II $\to$ P;III $\to$ R;IV $\to$ S}
\opt{C}{I $\to$ R;II $\to$ U;III $\to$ U;IV $\to$ Q}
\opt{D}{I $\to$ S;II $\to$ Q;III $\to$ R;IV $\to$ S}
\end{qoptions}
\end{bankq}
\begin{bankq}[id=TOAS05-P1-PHY-S2-Q10, source={TOAS - Test 05 (03 Apr 2022)}, paper=1, qno=10, type=single, subtopic={Spring-mass}, answer={B}, answer_src=key, difficulty=5, diff_parts={K1 S4 I1 C0 F0}, diff_note={Same events with phi=pi/4: each time found from the phase}]
A particle is performing SHM $x=A\sin(\omega t+\phi)$ whose time period is $T$. List-I contains events related to the motion of particle and List-II contains the minimum time after which that event happen. Choose the correct match from option if $\phi=\pi/4$
{\small\setlength{\tabcolsep}{3pt}\renewcommand{\arraystretch}{1.25}\begin{center}\begin{tabular}{|L{0.75cm}|L{3.6cm}|L{0.7cm}|L{1.0cm}|}
\hline
\multicolumn{2}{|c|}{List--I} & \multicolumn{2}{c|}{List--II}\\
\hline
(I) & Particle comes to rest & (P) & $T/4$\\
\hline
(II) & Velocity is maximum & (Q) & $T/6$\\
\hline
(III) & Kinetic energy = Potential energy & (R) & $3T/8$\\
\hline
(IV) & Velocity is half of maximum velocity & (S) & $T/24$\\
\hline
\multicolumn{2}{|c|}{} & (T) & $T/2$\\
\hline
\multicolumn{2}{|c|}{} & (U) & $T/8$\\
\hline
\end{tabular}\end{center}}
\begin{qoptions}
\opt{A}{I $\to$ P;II $\to$ T;III $\to$ U;IV $\to$ Q}
\opt{B}{I $\to$ U;II $\to$ R;III $\to$ P;IV $\to$ S}
\opt{C}{I $\to$ R;II $\to$ Q;III $\to$ S;IV $\to$ T}
\opt{D}{I $\to$ S;II $\to$ Q;III $\to$ R;IV $\to$ T}
\end{qoptions}
\end{bankq}
\begin{bankq}[id=TOAS05-P1-PHY-S3-Q02, source={TOAS - Test 05 (03 Apr 2022)}, paper=1, qno=2, type=numerical, subtopic={Spring-mass}, answer={2.66 to 2.67}, answer_src=key, difficulty=6, diff_parts={K2 S3 I1 C0 F1}, diff_note={Spring force to lift the 5 kg block: amplitude 1 cm + extension needed, from k=3g/0.01}]
The ends of a spring are attached to two blocks of masses 5 kg and 3 kg as shown in figure. The 5 kg block rests on a horizontal surface and the 3 kg block which is vertically above it is in equilibrium with the spring compressed by 1 cm. The 3 kg block is pushed down further by $x$ cm and released. Find the minimum value of $x$ so that the 5 kg block may be lifted off the ground during the subsequent motion.
\bfig{81.6}{TOAS05_p1_phy_s3_q02_1.png}
\end{bankq}
\begin{bankq}[id=TOAS05-P1-PHY-S3-Q06, source={TOAS - Test 05 (03 Apr 2022)}, paper=1, qno=6, type=numerical, subtopic={Spring-mass}, answer={9.00}, answer_src=key, difficulty=4, diff_parts={K1 S2 I1 C0 F1}, diff_note={8 cos^2(50t) = 4+4 cos 100t: amplitude 5 about x=4, max distance 9}]
A particle performs SHM along $x$-axis according to the relation $x=3\sin(100t)+8\cos^2(50t)$ where `$x$' is in cm and `$t$' is in second. Find the maximum possible distance (in cm) of the particle from origin during its motion.
\end{bankq}
\begin{bankq}[id=TOAS05-P2-PHY-S1-Q04, source={TOAS - Test 05 (03 Apr 2022)}, paper=2, qno=4, type=single, subtopic={Spring-mass}, answer={D}, answer_src=key, difficulty=2, diff_parts={K1 S2 I0 C0 F0}, diff_note={v=w sqrt(A^2-x^2) with w=4}]
A particle is performing SHM and its equation is $x=10\sin\left(4t-\dfrac\pi3\right)$, where time ($t$) is in second and $x$ is the position from the mean position (in metre). The speed of particle at the instant when its displacement from mean position is 6m, is:-
\begin{qoptions}
\opt{A}{$\dfrac{16\pi}3$ m/s}
\opt{B}{$\dfrac{10\pi}3$ m/s}
\opt{C}{18 m/s}
\opt{D}{32 m/s}
\end{qoptions}
\end{bankq}
\begin{bankq}[id=TOAS05-P2-PHY-S1-Q09, source={TOAS - Test 05 (03 Apr 2022)}, paper=2, qno=9, type=single, subtopic={Pendulums}, answer={B}, answer_src=key, difficulty=5, diff_parts={K2 S3 I1 C0 F0}, diff_note={Tension at extreme and mean positions gives the friction needed}]
Two bodies of mass $m$ and $4m$ are attached by a string shown in the figure. The body of mass $m$ hanging from a string of length $\ell$ is executing simple harmonic motion with small amplitude $A$ while other body is at rest on the surface. The minimum coefficient of friction between the mass $4m$ and the horizontal surface must be :
\bfig{137.9}{TOAS05_p2_phy_s1_q09_1.png}
\begin{qoptions}
\opt{A}{$\dfrac14\left(1-\dfrac{A^2}{\ell^2}\right)$}
\opt{B}{$\dfrac14\left(1+\dfrac{A^2}{\ell^2}\right)$}
\opt{C}{$\dfrac14\dfrac A\ell\cos\theta$}
\opt{D}{$\dfrac14$}
\end{qoptions}
\end{bankq}
\begin{bankq}[id=TOAS05-P2-PHY-S2-Q03, source={TOAS - Test 05 (03 Apr 2022)}, paper=2, qno=3, type=numerical, subtopic={Spring-mass}, answer={6}, answer_src=key, difficulty=5, diff_parts={K2 S3 I0 C0 F1}, diff_note={Elastic collision of identical blocks exchanges velocities; new amplitude v/w}]
One end of an ideal spring is fixed to a wall at origin $O$ and axis of the spring is parallel to $x$-axis. A block of mass 1 kg is attached to the free end of the spring and is performing SHM. Equation of position of the block in co-ordinate system shown in figure is $x=10+3\sin(10t)$, where $t$ is in second and $x$ in cm. Another identical block $B$, moving towards the origin with velocity 0.6 m/s collides elastically with the block $A$ at $t=0$. The new amplitude (in cm) of oscillations will be :
\bfig{131.1}{TOAS05_p2_phy_s2_q03_1.png}
\end{bankq}
\begin{bankq}[id=TOAS05-P2-PHY-S2-Q04, source={TOAS - Test 05 (03 Apr 2022)}, paper=2, qno=4, type=numerical, subtopic={Spring-mass}, answer={5}, answer_src=key, difficulty=5, diff_parts={K1 S3 I1 C0 F1}, diff_note={Both springs add: equilibrium extensions 20 and 10 cm give k_eff=2k and w^2=20 g}]
Two identical springs are attached to a small block $P$. The other ends of the springs are fixed at $A$ and $B$ such that $AB$ is vertical. When $P$ is in equilibrium the extension of top spring is 20 cm and extension of bottom spring is 10 cm. If the period of small vertical oscillations of $P$ about its equilibrium position is $\dfrac{\pi}{N\sqrt2}$ sec, then fill the value of $N$.
\bfig{63.0}{TOAS05_p2_phy_s2_q04_1.png}
\end{bankq}
\begin{bankq}[id=TOAS05-P2-PHY-S2-Q07, source={TOAS - Test 05 (03 Apr 2022)}, paper=2, qno=7, type=numerical, subtopic={Spring-mass}, answer={3}, answer_src=solved, working={Positions $\pm x_0$ with $|v|=v_{max}/4$ give $\cos\theta=\pm\frac14$; opposite velocities at opposite positions mean the motions are $\pi$ out of phase: $3\pi/x$ with $x=3$}, difficulty=5, diff_parts={K1 S3 I1 C0 F1}, diff_note={Speeds v_max/4 at +-x0 with opposite velocity: phase difference pi}]
Two particles are performing simple harmonic motion with same amplitude and same frequency. When they are at same distance from mean position on opposite sides, their speeds are one fourth of maximum speed. Also at these positions their direction of velocities are opposite. If the phase difference between the two simple harmonic motions is $\dfrac{3\pi}{x}$, then find the value of $x$.
\end{bankq}
\begin{bankq}[id=TOAS05-P2-PHY-S2-Q10, source={TOAS - Test 05 (03 Apr 2022)}, paper=2, qno=10, type=numerical, subtopic={Pendulums}, answer={4}, answer_src=key, difficulty=6, diff_parts={K2 S3 I1 C0 F1}, diff_note={Physical pendulum: I=6mR^2 about the pivot, centre of mass at 1.5R}]
A ring of mass $m$ is pivoted at a point on its periphery and free to oscillates in vertical plane. A particle of same mass is attached at the lowest point of the ring then time period (in second) of small oscillations is [use $R=2$m ; $g=\pi^2$ m/s$^2$]
\bfig{85.8}{TOAS05_p2_phy_s2_q10_1.png}
\end{bankq}
\begin{bankq}[id=TOAS06-P1-PHY-S1-Q04, source={TOAS - Test 06 (17 Apr 2022)}, paper=1, qno=4, type=single, subtopic={Pendulums}, answer={C}, answer_src=key, difficulty=6, diff_parts={K2 S3 I1 C1 F0}, diff_note={Bar on circular supports: centre of mass height from the geometry, energy for small angle gives the frequency}]
The bar has a length $\ell$ and mass $m$. It is supported at its ends by rollers of negligible mass. If it is given a small displacement and released, determine the natural frequency of vibration.
\bfig{128.5}{TOAS06_p1_phy_s1_q04_B.png}
\begin{qoptions}
\opt{A}{$\omega_N=\sqrt{\dfrac{3g\left(4R^{2}-\ell^{2}\right)}{6R^{2}-\ell^{2}}}$}
\opt{B}{$\omega_N=\sqrt{\dfrac{3g\left(4R^{2}-\ell^{2}\right)^{\frac12}}{6R^{2}+\ell^{2}}}$}
\opt{C}{$\omega_N=\sqrt{\dfrac{3g\left(4R^{2}-\ell^{2}\right)^{\frac12}}{6R^{2}-\ell^{2}}}$}
\opt{D}{$\omega_N=\sqrt{\dfrac{3g\left(4R^{2}+\ell^{2}\right)^{\frac12}}{6R^{2}-\ell^{2}}}$}
\end{qoptions}
\end{bankq}
\begin{bankq}[id=TOAS07-P1-PHY-S1-Q03, source={TOAS - Test 07 (01 May 2022)}, paper=1, qno=3, type=single, subtopic={Spring-mass}, answer={D}, answer_src=key, difficulty=5, diff_parts={K2 S3 I1 C0 F0}, diff_note={Speed at mean position, perfectly inelastic sticking, new amplitude reaches the wall}]
A small block of mass $m$ resting on a horizontal smooth floor at distance $\ell$ from a rigid wall. Block attached to the spring, is pushed towards right by a distance $\dfrac{3\ell}{2}$ and released. When block passes from its mean position an other block of mass $m_1$ is placed on it which sticks to it due to friction. If combined block just collides with the left wall, then :-
\bfig{114.4}{TOAS07_p1_phy_s1_q03_1.png}
\begin{qoptions}
\opt{A}{$m_1=m$}
\opt{B}{$m_1=\dfrac m4$}
\opt{C}{$m_1=\dfrac{3m}{4}$}
\opt{D}{$m_1=\dfrac{5m}{4}$}
\end{qoptions}
\end{bankq}
\begin{bankq}[id=TOAS07-P1-PHY-S2-Q09, source={TOAS - Test 07 (01 May 2022)}, paper=1, qno=9, type=multi, subtopic={Spring-mass}, answer={A,D}, answer_src=key, difficulty=5, diff_parts={K1 S3 I1 C0 F1}, diff_note={Amplitude from energy/velocity at mean position, period from added mass}]
A horizontal spring mass system of mass $M$ executes oscillatory motion of amplitude $a_0$ and time period $T_0$. When the mass $M$ is passing through its equilibrium position another mass $m$ is placed on it such that both move together. If $a$ and $T$ be the new amplitude and time period respectively then,
\begin{qoptions}
\opt{A}{$a=\sqrt{\dfrac{M}{M+m}}\,a_0$}
\opt{B}{$a=\sqrt{\dfrac{M+m}{M}}\,a_0$}
\opt{C}{$T=\sqrt{\dfrac{M}{M+m}}\,T_0$}
\opt{D}{$T=\sqrt{\dfrac{M+m}{M}}\,T_0$}
\end{qoptions}
\end{bankq}
\begin{bankq}[id=TOAS07-P2-PHY-S1-Q02, source={TOAS - Test 07 (01 May 2022)}, paper=2, qno=2, type=multi, subtopic={Spring-mass}, answer={A,B}, answer_src=key, difficulty=6, diff_parts={K2 S3 I1 C0 F1}, diff_note={Test charge between plates with two springs: SHM about displaced centre}]
Two large non conducting plates having surface charge densities $+\sigma$ and $-\sigma$, respectively, are fixed `$d$' distance apart. A small test charge $q$ of mass $m$ is attached to two non-conducting identical springs of spring constant $k$ as shown in the adjacent figure. The charge $q$ is now released from rest with springs in natural length. Then $q$ will [neglect gravity] :-
\bfig{115.9}{TOAS07_p2_phy_s1_q02_1.png}
\begin{qoptions}
\opt{A}{perform SHM with angular frequency $\sqrt{\dfrac{2k}{m}}$}
\opt{B}{perform SHM with amplitude $\dfrac{\sigma q}{2k\epsilon_0}$}
\opt{C}{not perform SHM, but will have a periodic motion if charges are removed on plates as well as on $m$}
\opt{D}{remain stationary}
\end{qoptions}
\end{bankq}
\begin{bankq}[id=TOAS07-P2-PHY-S1-Q04, source={TOAS - Test 07 (01 May 2022)}, paper=2, qno=4, type=multi, subtopic={Spring-mass}, answer={B,C}, answer_src=key, difficulty=7, diff_parts={K2 S3 I1 C1 F1}, diff_note={Pulley-spring system: energy for max speed and maximum elongation}]
Block $m$ is released from rest when spring is in its natural length (assume pulley is ideal and block does not strike on ground during it's motion in vertical plane) then :
\bfig{54.0}{TOAS07_p2_phy_s1_q04_1.png}
\begin{qoptions}
\opt{A}{maximum speed of block is $g\sqrt{\dfrac mk}$}
\opt{B}{maximum speed of block is $2g\sqrt{\dfrac mk}$}
\opt{C}{maximum elongation in spring is $4\,mg/k$}
\opt{D}{maximum elongation in spring is $2\,mg/k$}
\end{qoptions}
\end{bankq}
\begin{bankq}[id=TOAS09-P1-PHY-S1-Q03, source={TOAS - Test 09 (30 May 2022)}, paper=1, qno=3, type=single, subtopic={Spring-mass}, answer={C}, answer_src=key, difficulty=4, diff_parts={K2 S2 I1 C0 F0}, diff_note={Spring-mass: standard approach for this subtopic; 2 concept(s), 2 step(s), needs a non-obvious insight}]
All the springs and pulleys are ideal and massless. If the block is slightly pulled vertically downward from its equilibrium position, then the time period of oscillation of the block is (assume strings are taught all the time)
\bfig{118.0}{TOAS09_p1_phy_s1_q03_B.png}
\begin{qoptions}
\opt{A}{$2\pi\sqrt{\dfrac{m}{4k}}$}
\opt{B}{$\pi\sqrt{\dfrac{2m}{k}}$}
\opt{C}{$\pi\sqrt{\dfrac{m}{2k}}$}
\opt{D}{$2\pi\sqrt{\dfrac{m}{16k}}$}
\end{qoptions}
\end{bankq}
\begin{bankq}[id=TOAS09-P2-PHY-S1-Q17, source={TOAS - Test 09 (30 May 2022)}, paper=2, qno=17, type=single, subtopic={Pendulums}, answer={C}, answer_src=key, difficulty=4, diff_parts={K2 S2 I1 C0 F0}, diff_note={Pendulums: standard approach for this subtopic; 2 concept(s), 2 step(s), needs a non-obvious insight}]
Time period of a simple pendulum whose length is $\dfrac R2$, where $R$ is radius of the earth oscillating near earth's surface will be -
\begin{qoptions}
\opt{A}{$2\pi\sqrt{\dfrac{R}{2g}}$}
\opt{B}{$2\pi\sqrt{\dfrac{R}{g}}$}
\opt{C}{$2\pi\sqrt{\dfrac{R}{3g}}$}
\opt{D}{$2\pi\sqrt{\dfrac{3R}{g}}$}
\end{qoptions}
\end{bankq}
\end{document}
